\clearpage\pagestyle{otMain}\setlength{\otBlocEspacement}{2.5mm}\fontsize{11}{15}\selectfont\corrigesChapitre{9}{Relier des grandeurs et lire des graphiques}
\begin{blocLivre}{cor-9-1}\begin{corrige}{1}{Une formule de tarif}\demarche{Additionner partie fixe et partie variable.}
La partie fixe vaut 5 et la partie liée au temps vaut $3t$. Donc $\res{P=5+3t}$, avec $P$ en euros et $t$ en heures.
\end{corrige}\end{blocLivre}
\begin{blocLivre}{cor-9-2}\begin{corrige}{2}{Compléter les valeurs}\demarche{Substituer chaque durée dans la même règle.}
On remplace $t$ dans la formule : pour 0 heure, le prix vaut 5 euros ; les deux autres prix valent $5+6=11$ et $5+12=17$. On retient $\res{5,\ 11,\ 17}$ euros, sans avoir besoin d’une notation de fonction.
\end{corrige}\end{blocLivre}
\begin{blocLivre}{cor-9-3}\begin{corrige}{3}{Lire un couple}\demarche{Suivre la verticale puis l’horizontale.}
On part de 2 sur l’axe du temps, on rejoint le tracé puis on lit l’axe vertical : la température est $\res{60\,^\circ\mathrm C}$. Les unités viennent des axes.
\label{fig-cor-9-3}\figureGeo[corrige]{four}
\end{corrige}\end{blocLivre}
\begin{blocLivre}{cor-9-4}\begin{corrige}{4}{Lire une durée}\demarche{Partir cette fois de l’axe vertical.}
L’horizontale de niveau 100 coupe le tracé à l’abscisse $\res4$ min pour la première fois. On lit un temps, et non une température.
\label{fig-cor-9-4}\figureGeo[corrige]{four}
\end{corrige}\end{blocLivre}
\begin{blocLivre}{cor-9-5}\begin{corrige}{5}{Reconnaître une relation}\demarche{Vérifier le passage par l’origine.}
La droite représentée passe par $(0;2)$, pas par l’origine. La relation n’est donc pas proportionnelle, même si les points sont alignés. Pour zéro unité horizontale, la valeur verticale est déjà 2.
\label{fig-cor-9-5}\figureGeo[corrige]{graph-non-proportion}
\end{corrige}\end{blocLivre}
\begin{blocLivre}{cor-9-6}\begin{corrige}{6}{Construire un graphique}\demarche{Placer les couples avant de relier.}
Les prix sont $\res{0,\ 3,\ 6,\ 9}$ euros. Les points $(0;0)$, $(2;3)$, $(4;6)$ et $(6;9)$ s’alignent avec l’origine. La corde pouvant être achetée en longueur intermédiaire dans ce modèle, on trace le segment correspondant.
\label{fig-cor-9-6}\figureGeo[corrige]{graph-corde}
\end{corrige}\end{blocLivre}
\begin{blocLivre}{cor-9-7}\begin{corrige}{7}{Une lecture inverse double}\demarche{Parcourir toute l’horizontale.}
La droite horizontale de niveau 60 coupe le tracé deux fois : à $\res2$ min pendant la chauffe et à $\res{10}$ min pendant le refroidissement. Il faut chercher toutes les intersections dans l’intervalle représenté.
\label{fig-cor-9-7}\figureGeo[corrige]{four}
\end{corrige}\end{blocLivre}
\begin{blocLivre}{cor-9-8}\begin{corrige}{8}{Aire d’un carré}\demarche{Comparer un doublement de la grandeur de départ.}
$A=c^2$. Les aires valent $\res{1,4,9,16}$ cm$^2$. Doubler le côté de 1 à 2 multiplie l’aire par 4 et non par 2 : l’aire n’est pas proportionnelle au côté.
\label{fig-cor-9-8}\figureGeo[corrige]{graph-carre}
\end{corrige}\end{blocLivre}
\begin{blocLivre}{cor-9-9}\begin{corrige}{9}{Une formule à retrouver}\demarche{Utiliser la valeur pour zéro et la règle annoncée.}
Pour une première valeur $x$, le montant est $3x+b$. Le couple $(0;4)$ donne $b=4$. La formule est $\res{P=3x+4}$. Elle retrouve les trois autres couples. Sans l’hypothèse sur le modèle, le tableau seul ne fixerait pas une règle unique.
\end{corrige}\end{blocLivre}
\begin{blocLivre}{cor-9-10}\begin{corrige}{10}{Un tableur}\demarche{Remplacer la lettre par la cellule de la ligne.}
On saisit $\res{\texttt{=5+3*A2}}$. Lors de la recopie, A2 devient A3, A4, etc. La référence relative suit la durée de chaque ligne ; les nombres 5 et 3 restent fixes.
\label{fig-cor-9-10}\figureGeo[corrige]{tableur-prix}
\end{corrige}\end{blocLivre}
\begin{blocLivre}{cor-9-11}\begin{corrige}{11}{Retirer des carrés}\demarche{Soustraire les quatre aires retirées.}
L’aire initiale vaut 180 cm$^2$. Les quatre carrés retirent $4x^2$, donc $\res{A=180-4x^2}$. Pour $x=2$, $A=180-16=\res{164\ \mathrm{cm}^2}$. La condition $x<5$ empêche les carrés des deux bords courts de se chevaucher.
\label{fig-cor-9-11}\figureGeo[corrige]{coins-retires}
\end{corrige}\end{blocLivre}
\begin{blocLivre}{cor-9-12}\begin{corrige}{12}{Deux courbes de prix}\demarche{Confirmer la lecture par le calcul.}
Les tracés se coupent au point $(2;8)$. L’égalité $4+2t=4t$ donne $4=2t$, donc $\res{t=2}$ h et un prix de $\res8$ euros. Avant 2 h, B est moins cher ; après 2 h, A est moins cher sur l’intervalle étudié.
\label{fig-cor-9-12}\figureGeo[corrige]{graph-tarifs}
\end{corrige}\end{blocLivre}
\begin{blocLivre}{cor-9-13}\begin{corrige}{13}{Une réserve qui diminue}\demarche{Associer une formule à un domaine physique.}
Le volume restant est $\res{V=120-8t}$ L. Il atteint zéro pour $120-8t=0$, soit $t=15$ min. Le modèle de vidange est valable pour $\res{0\le t\le15}$ ; un volume négatif ne représente pas la situation.
\label{fig-cor-9-13}\figureGeo[corrige]{graph-reserve}
\end{corrige}\end{blocLivre}
\begin{blocLivre}{cor-9-14}\begin{corrige}{14}{Connaître les limites}\demarche{Distinguer données observées et prolongement supposé.}
Non. Une tendance sur l’intervalle janvier-juin ne fixe pas la suite. Il manque les données ou un modèle justifié pour décembre. Même une hausse régulière dans la zone observée ne prouve pas qu’elle continuera.
\end{corrige}\end{blocLivre}
\begin{blocLivre}{cor-9-15}\begin{corrige}{15}{Comparer deux dépendances}\demarche{Utiliser la positivité pour diviser par le côté.}
Pour $c=4$, $P=16$ et $A=16$. Si $c>0$, l’égalité numérique $c\times c=4c$ peut être divisée par $c$, non nul : on obtient $\res{c=4}$ seulement. Les unités restent différentes, cm et cm$^2$ ; on compare des valeurs numériques.
\label{fig-cor-9-15}\figureGeo[corrige]{carre-4}
\end{corrige}\end{blocLivre}
